\SourcePage{322}
\SGAsection{6}{局部消失圈理论的应用（Picard--Lefschetz）}
\begin{SGAresult}{命题}{6.1}
设 $Z$ 是\Term{projective-space}{射影空间} $\mathbb P_k^r$ 中的一个\Term{hypersurface}{超曲面}，其中 $r\geq2$，$k$ 是\Term{algebraically-closed-field}{代数闭域}。设
\[
i_0:D_0\hookrightarrow\mathbb P^r
\]
是一条直线的\Term{inclusion}{包含映射}；该直线与 $Z$ 仅在后者的\Term{smooth}{光滑}\Term{open-subset}{开集}中相交，并且在那里与之\Term{transverse-intersection}{横截相交}。选取 $D_0-D_0\cap Z$ 的一个\Term{geometric-point}{几何点} $\bar a$，并用
\[
\pi_1^{\mathrm{mod}}(D_0-D_0\cap Z,\bar a)
\]
表示 $\pi_1(D_0-D_0\cap Z,\bar a)$ 的如下\Term{quotient-group}{商群}：它分类 $D_0-D_0\cap Z$ 上在 $D_0\cap Z$ 每一点均\Term{tamely-ramified}{温和分歧}的\Term{etale}{平展}\Term{covering}{覆盖}。则：
\SourcePage{323}
\textbf{6.1.1.} \Term{canonical}{典范}\Term{homomorphism}{同态}
\[
i_{0*}:\pi_1^{\mathrm{mod}}(D_0-D_0\cap Z,\bar a)
\longrightarrow\pi_1^{\mathrm{mod}}(\mathbb P^r-Z,\bar a)
\tag{6.1.1.1}
\]
是\Term{surjective}{满射}。

\textbf{6.1.2.} 每选取一个\Term{geometric-point}{几何点} $\bar x_0\in D_0\cap Z$ 以及 $D_0$ 中一条连接 $\bar x_0$ 与 $\bar a$ 的“路径”，便确定一个\Term{homomorphism}{同态}
\[
\prod_{\ell\neq\operatorname{car}(k)}\mathbb Z_\ell(1)
\longrightarrow\pi_1^{\mathrm{mod}}(D_0-D_0\cap Z,\bar a),
\tag{6.1.2.1}
\]
其像就是 $\bar x_0$ 处的“那个”\Term{inertia-group}{惯性群}。将此\Term{morphism}{态射}与 $i_{0*}$ 复合，得到\Term{homomorphism}{同态}
\[
\prod_{\ell\neq\operatorname{car}(k)}\mathbb Z_\ell(1)
\longrightarrow\pi_1^{\mathrm{mod}}(\mathbb P^r-Z).
\tag{6.1.2.2}
\]
若 $Z$ \Term{irreducible}{不可约}，则所有这些\Term{homomorphism}{同态}彼此\Term{conjugacy}{共轭}。
\end{SGAresult}
\begin{proof}
用 $T$ 表示 $\mathbb P^r$ 中直线的 Grassmann 簇的如下\Term{subscheme}{子概形}，其参数化的直线满足：
\begin{enumerate}
\item[(i)] 通过 $\bar a$；
\item[(ii)] 与 $Z$ 仅在后者的\Term{smooth}{光滑}\Term{open-subset}{开集}中相交，并且在那里与之\Term{transverse-intersection}{横截相交}。
\end{enumerate}
用 $\eta$ 表示 $T$ 的\Term{generic-point}{泛点}，用 $i_\eta:D_\eta\hookrightarrow\mathbb P^r\otimes_k k(\eta)$ 表示相应的直线（因而它是一条“将 $\bar a$ 与 $\mathbb P_k^r$ 的\Term{generic-point}{泛点}连接起来”的直线）。再选取 $\eta$ 上方的一个\Term{geometric-point}{几何点} $\bar\eta$，并用 $D_{\bar\eta}$ 表示相应的“几何”直线。令 $o$ 为 $T$ 中对应于给定直线 $D_0$ 的点。

\SourcePage{324}
由 SGA 1，第 XIII 讲 2.8，存在一个\Term{specialization-homomorphism}{特殊化同态}（它依赖于 $T$ 中一条连接 $\bar\eta$ 与 $o$ 的路径，故只确定到\Term{conjugacy}{共轭}）
\[
\pi_1^{\mathrm{mod}}(D_{\bar\eta}-D_{\bar\eta}\cap Z,\bar a)
\xrightarrow{\bar\eta\to o}
\pi_1^{\mathrm{mod}}(D_0-D_0\cap Z,\bar a)
\tag{6.1.3}
\]
并满足：

\textbf{6.1.4.} \Term{local-scheme}{局部}\Term{monodromy}{单值性}\Term{homomorphism}{同态} (6.1.2.1)
\[
\prod_{\ell\neq\operatorname{car}(k)}\mathbb Z_\ell(1)
\longrightarrow\pi_1(D_0-D_0\cap Z,\bar a)
\]
可由 $D_{\bar\eta}$ 的\Term{local-scheme}{局部}\Term{monodromy}{单值性}\Term{homomorphism}{同态}再与\Term{specialization-homomorphism}{特殊化同态} (6.1.3) 复合而得到：
\[
\prod_{\ell\neq\operatorname{car}(k)}\mathbb Z_\ell(1)
\longrightarrow\pi_1^{\mathrm{mod}}(D_{\bar\eta}-D_{\bar\eta}\cap Z,\bar a)
\xrightarrow{\bar\eta\to o}\pi_1^{\mathrm{mod}}(D_0-D_0\cap Z,\bar a).
\tag{6.1.4.1}
\]

\textbf{6.1.5.} 下图交换：
\[
\begin{tikzcd}[column sep=small,row sep=large]
&\pi_1^{\mathrm{mod}}(\mathbb P^r-Z,\bar a)&\\
\pi_1^{\mathrm{mod}}(D_\eta-D_\eta\cap Z,\bar a)
\arrow[ur,"i_{\eta*}"]&&
\pi_1^{\mathrm{mod}}(D_0-D_0\cap Z,\bar a)\arrow[ul,"i_{0*}"']\\
&\pi_1^{\mathrm{mod}}(D_{\bar\eta}-D_{\bar\eta}\cap Z,\bar a)
\arrow[ul]\arrow[uu,"i_{\bar\eta*}" near start]
\arrow[ur,"\bar\eta\to o"']&
\end{tikzcd}
\tag{6.1.5}
\]

\SourcePage{325}
先承认上述事实，并证明 6.1.1。由 (6.1.5) 的\Term{commutativity}{交换性}，只需证明\Term{homomorphism}{同态}
\[
i_{\bar\eta*}:\pi_1^{\mathrm{mod}}(D_{\bar\eta}-D_{\bar\eta}\cap Z,\bar a)
\longrightarrow\pi_1^{\mathrm{mod}}(\mathbb P^r-Z,\bar a)
\]
是\Term{surjective}{满射}；为此，只需证明不取温和商的\Term{fundamental-group}{基本群}之间的相应\Term{homomorphism}{同态}
\[
i_{\bar\eta}:\pi_1(D_{\bar\eta}-D_{\bar\eta}\cap Z,\bar a)
\longrightarrow\pi_1(\mathbb P^r-Z,\bar a)
\tag{6.1.6}
\]
是\Term{surjective}{满射}。也就是说，对每个\Term{etale}{平展}\Term{covering}{覆盖} $E\to\mathbb P^r-Z$，若 $E$ \Term{connected}{连通}（从而\Term{irreducible}{不可约}，因为 $\mathbb P^r-Z$ \Term{smooth}{光滑}，因而\Term{unibranch}{单支}，见 EGA IV 18.10.3），则它在 $D_{\bar\eta}-D_{\bar\eta}\cap Z$ 上的\Term{inverse-image}{逆像}
\[
E_{\bar\eta}=E\mathbin{\times}_{\mathbb P^r-Z}
(D_{\bar\eta}-D_{\bar\eta}\cap Z)
\]
仍然\Term{geometrically-irreducible}{几何不可约}。为此，只需使用以下 Bertini \Term{theorem}{定理}的变式（EGA V \SGAsourcegap{原文此处引注仅作省略点，未给出 EGA V 中的具体位置}），将它用于复合\Term{arrow}{箭头} $E\to\mathbb P^r-Z\hookrightarrow\mathbb P^r$，并取 $s=1$：

\begin{SGAresult}{Bertini \Term{theorem}{定理}}{6.1.6.1}
设 $k$ 是一个域，$\mathbb P^r=\mathbb P_k^r$，$f:X\to\mathbb P^r$ 是一个非常值\Term{morphism}{态射}，其中 $X$ 是一个\Term{geometrically-irreducible}{几何不可约}的\Term{algebraic-scheme}{代数概形}。对每个整数 $s<\dim\overline{f(X)}$，用 $L_t$ 表示 $\mathbb P^r$ 中\Term{codimension}{余维数}为 $s$ 的泛\Term{linear-subspace}{线性子空间}。则 $X_t=f^{-1}(L_t)$ 是\Term{geometrically-irreducible}{几何不可约}的。
\end{SGAresult}
现在证明 6.1.2。由 6.1.4 以及 (6.1.5) 的\Term{commutativity}{交换性}，只需证明 $D_{\bar\eta}$ 的所有\Term{local-scheme}{局部}\Term{monodromy}{单值性}\Term{homomorphism}{同态} (6.1.3)
\SourcePage{326}
\[
\prod_{\ell\neq\operatorname{car}(k)}\mathbb Z_\ell(1)
\longrightarrow\pi_1^{\mathrm{mod}}(D_{\bar\eta}-D_{\bar\eta}\cap Z,\bar a)
\tag{6.1.7.1}
\]
在与\Term{homomorphism}{同态}
\[
\pi_1^{\mathrm{mod}}(D_{\bar\eta}-D_{\bar\eta}\cap Z,\bar a)
\longrightarrow\pi_1^{\mathrm{mod}}(D_\eta-D_\eta\cap Z,\bar a)
\tag{6.1.7.2}
\]
复合之后彼此\Term{conjugacy}{共轭}。

为此，注意到，对任意\Term{finite}{有限}\Term{galois-extension}{Galois 扩张} $K/k(\eta)$，其中 $k(\eta)\subset K\subset k(\bar\eta)$，\Term{scheme}{概形}
\[
D_K\overset{\mathrm{def}}=D_\eta\times_{k(\eta)}K
\]
是 $D_\eta$ 的一个\Term{etale}{平展}\Term{covering}{覆盖}，所以 $D_K-D_K\cap Z$ 是 $D_\eta-D_\eta\cap Z$ 的一个\Term{tamely-ramified}{温和}\Term{covering}{覆盖}，由此得到 $\pi_1^{\mathrm{mod}}(D_\eta-D_\eta\cap Z,\bar a)$ 在\Term{scheme}{概形} $D_K\cap Z$ 上的\Term{action}{作用}。取 $K$ 充分大，使得 $D_K\cap Z$ 的所有点均为 $K$ \Term{rational-point}{有理点}；这是可能的，因为由假设，\Term{scheme}{概形} $D_\eta\cap Z$ 在 $k(\eta)$ 上\Term{etale}{平展}。

为证各个 $D_K\cap Z$ 中的点所对应的\Term{local-scheme}{局部}\Term{inertia-group}{惯性群}的复合\Term{homomorphism}{同态}
\[
\prod_{\ell\neq\operatorname{car}(k)}\mathbb Z_\ell(1)
\longrightarrow\pi_1^{\mathrm{mod}}(D_\eta-D_\eta\cap Z,\bar a)
\]
彼此\Term{conjugacy}{共轭}，只需证明群 $\pi_1^{\mathrm{mod}}(D_\eta-D_\eta\cap Z,\bar a)$ 在 $D_K\cap Z$ 的点集上\Term{transitive-action}{传递地作用}。而 $\pi_1^{\mathrm{mod}}(D_\eta-D_\eta\cap Z,\bar a)$ 在 $D_K\cap Z$ 上的\Term{action}{作用}通过\Term{canonical}{典范}商
\[
\pi_1^{\mathrm{mod}}(D_\eta-D_\eta\cap Z,\bar a)
\longrightarrow\pi_1(\eta,\bar\eta)=\operatorname{Gal}(/k(\eta))
\]
\Term{resolution}{分解}，
\SourcePage{327}
且 $\operatorname{Gal}(/k(\eta))$ 在 $D_K\cap Z$ 上的\Term{action}{作用}与\Term{inclusion}{包含映射}
\[
D_K\cap Z\hookrightarrow\mathbb P_k^r\times_kK
\]
相容，其中 $\operatorname{Gal}(/k(\eta))$ 只通过第二个因子\Term{action}{作用}于 $\mathbb P_k^r\times_kK$。于是，所要的结论\Term{equivalence}{等价}于\Term{scheme}{概形} $D_\eta\cap Z$ 只有一个点（即它是一个域的谱）。用最直接的语言表述，这恰好就是如下 Bertini 型事实（为 $Z$ 取一个\Term{affine}{仿射}方程 $F=0$，并设 $\bar a$ 是原点）：

\begin{SGAresult}{事实}{6.1.8}
设 $k$ 是\Term{algebraically-closed-field}{代数闭域}，$r\geq2$，$F(X_1,\ldots,X_r)\in k[X_1,\ldots,X_r]$ 是一个\Term{irreducible}{不可约}\Term{polynomial}{多项式}。则\Term{polynomial}{多项式}
\[
G(T)=F(X_1T,\ldots,X_rT)\in k(X_1,\ldots,X_r)[T]
\]
是\Term{irreducible}{不可约}的（视为系数在 $k(X_1,\ldots,X_r)$ 中的关于 $T$ 的\Term{polynomial}{多项式}）。
\end{SGAresult}
这就完成了 6.1.2 的证明。
\end{proof}

\textbf{注记 6.1.9。} 另一种证明 6.1.2 的方法在明显更\Term{generic}{一般}的条件下仍然\Term{effective}{有效}：只需确认“\Term{monodromy}{单值性}\Term{homomorphism}{同态}”的\Term{generic}{一般}相容性，即 6.1.2 中的各个复合\Term{homomorphism}{同态}都与 $Z$ 的\Term{generic-point}{泛点}所对应的\Term{monodromy}{单值性}\Term{homomorphism}{同态}\Term{conjugacy}{共轭}（从而它们彼此\Term{conjugacy}{共轭}！）。

\SourcePage{328}
\begin{SGAresult}{推论}{6.1.10}
设 $X$ 是\Term{algebraically-closed-field}{代数闭域} $k$ 上的 $n$ 维\Term{smooth}{光滑}\Term{connected}{连通}\Term{projective}{射影}\Term{algebraic-variety}{簇}。设 $n=\dim X$ 为偶数，或者 $\operatorname{card}(k)\neq2$。设 $D$ 是 $X$ 的一个\Term{lefschetz-pencil}{Lefschetz 铅笔}，$\eta$ 是它的\Term{generic-point}{泛点}，$\bar\eta$ 是 $\eta$ 上方的一个\Term{geometric-point}{几何点}。则 $\pi_1(D-D\cap\check X,\bar\eta)$ 在 $H^{n-1}(X_{\bar\eta},\mathbb Q_\ell)$ 上的自然\Term{representation}{表示}通过 $\pi_1^{\mathrm{mod}}(D-D\cap\check X,\bar\eta)$ \Term{resolution}{分解}。若此\Term{representation}{表示}非平凡，则 $D\cap\check X\neq\varnothing$，且 $D\cap\check X$ 每一点的\Term{inertia-group}{惯性群}都非平凡地\Term{action}{作用}。
\end{SGAresult}
\begin{proof}
若 $\dim(\check X)\leq r-2$，则 $D\cap\check X$ 为空；考虑到 $\mathbb P^1$ \Term{simply-connected}{单连通}，此时没有什么需要证明。否则 $\check X$ 是 $\mathbb P^r$ 中的一个\Term{hypersurface}{超曲面}（第 XVII 讲 3.1.4），并且由 3.2.10 及第 XVII 讲 3.5，可以应用 6.1.1（取 $\check X$ 为 $Z$，取 $D$ 为 $D_0$）。$\pi_1$ 在 $H^{n-1}(X_{\bar\eta},\mathbb Q_\ell)$ 上的\Term{action}{作用}是\Term{tame-action}{温和作用}，这一点已在\Term{local-scheme}{局部}理论中借助\Term{picard-lefschetz-formula}{Picard--Lefschetz 公式}（第 XV 讲 3.4）看出；下文还将回顾该公式。再考虑到如下并非平凡的事实：$\pi_1^{\mathrm{mod}}(D-D\cap\check X,\bar\eta)$ 由 $D\cap\check X$ 各点的\Term{inertia-group}{惯性群}生成（SGA 1，第 XIII 讲 2.10），即得所需结论。
\end{proof}
现在用上述结果研究 (5.3.5) 的条件 (A)。先回顾以下结果（第 XV 讲 3.4），它是\Term{picard-lefschetz-formula}{Picard--Lefschetz 公式}的一个简单推论：

\SourcePage{329}
\begin{SGAresult}{命题}{6.2}
设 $X$ 是\Term{algebraically-closed-field}{代数闭域} $k$ 上的 $n$ 维\Term{smooth}{光滑}\Term{connected}{连通}\Term{projective}{射影}\Term{algebraic-variety}{簇}。用 $D$ 表示一个\Term{lefschetz-pencil}{Lefschetz 铅笔}，$\eta$ 表示它的\Term{generic-point}{泛点}，$\bar\eta$ 表示 $\eta$ 上方的一个\Term{geometric-point}{几何点}。则\Term{pencil}{铅笔} $D$ 满足条件 (A)，当且仅当 $D\cap\check X$ 每一点的\Term{inertia-group}{惯性群}都在 $H^{n-1}(X_{\bar\eta},\mathbb Q_\ell)$ 上非平凡地\Term{action}{作用}。
\end{SGAresult}

\begin{SGAresult}{\Term{theorem}{定理}}{6.3}
在 6.2 的假设下，有：

\textbf{6.3.1.} 若 $n=\dim X$ 为奇数，则 $X$ 的每个\Term{lefschetz-pencil}{Lefschetz 铅笔}都满足条件 (A)。

\textbf{6.3.2.} 设 $n=\dim X$ 为偶数，或者 $\operatorname{car}(k)\neq2$。则下列条件\Term{equivalence}{等价}：
\begin{enumerate}
\item[(i)] $D$ 不满足条件 (A)（5.3.5）；
\item[(ii)] $\pi_1(D-D\cap\check X)$ 在 $H^{n-1}(X_{\bar\eta},\mathbb Q_\ell)$ 上平凡地\Term{action}{作用}，且 $D\cap\check X\neq\varnothing$；
\item[(iii)] 存在 $D\cap\check X$ 中的一个点，其\Term{inertia-group}{惯性群}在 $H^{n-1}(X_{\bar\eta},\mathbb Q_\ell)$ 上平凡地\Term{action}{作用}；
\item[(iv)] \Term{sheaf}{层} $R^{n-1}p_*\mathbb Q_\ell$ 是 $D$ 上的\Term{constant-sheaf}{常值层}，且 $D\cap\check X$ \Term{nonempty}{非空}；
\item[(v)] \Term{sheaf}{层} $R^{n-1}p_*\mathbb Q_\ell|U$ 是 $U$ 上的\Term{constant-sheaf}{常值层}（其中 $U$ 是 $D$ 的一个\Term{nonempty}{非空}\Term{open-subset}{开集}，且 $p$ 在 $U$ 上\Term{smooth}{光滑}），并且 $p$ 不\Term{smooth}{光滑}。
\end{enumerate}
\SourcePage{330}
\textbf{6.3.3.} 若 $X$ 还满足条件 (LV)（5.2.2），则上述条件 (i)--(v) 还\Term{equivalence}{等价}于：
\begin{enumerate}
\item[(vi)] \Term{restriction-map}{限制映射}
\[
H^{n-1}(X,\mathbb Q_\ell)\longrightarrow H^{n-1}(X_{\bar\eta},\mathbb Q_\ell)
\]
是\Term{surjective}{满射}（因而由“\Term{weak-lefschetz-theorem}{弱 Lefschetz 定理}”，它是一个\Term{isomorphism}{同构}）。
\end{enumerate}
\end{SGAresult}
\begin{proof}
断言 6.3.1 由第 XV 讲 3.4 得出。

再看 6.3.2 中的\Term{equivalence}{等价}性：由 6.2 有 $(\mathrm{i})\Longleftrightarrow(\mathrm{iii})$，由 6.1.6 有 $(\mathrm{ii})\Longleftrightarrow(\mathrm{iii})$，$(\mathrm{ii})\Longleftrightarrow(\mathrm{v})$ 由定义即明，且由 5.3.4 有 $(\mathrm{iv})\Longleftrightarrow(\mathrm{v})$。最后，若 $X$ 满足条件 (LV)，则 5.6.10 给出 $(\mathrm{ii})\Longleftrightarrow(\mathrm{vi})$。
\end{proof}

\textbf{注记 6.3.4。} 可见，条件 (A) 不依赖于\Term{lefschetz-pencil}{Lefschetz 铅笔} $D$ 的选择，而只依赖于所考虑的（\Term{lefschetz-embedding}{Lefschetz}）\Term{immersion}{浸入} $i:X\to\mathbb P^r$。事实上，若 $\dim X$ 为奇数，则由 (6.3.1) 这是显然的；若 $\dim X$ 为偶数，则条件 (A) 也表示 $\mathbb P^r-\check X$ 上由各个 $H^{n-1}(X\cap H_s,\mathbb Q_\ell)$ 构成的\Term{local-system}{局部系统}是常值的。这由判据 (ii)、所考虑的\Term{local-system}{局部系统}在 $\check X$ 的\Term{generic-point}{泛点}处\Term{tamely-ramified}{温和分歧}这一事实，以及 (6.1.2) 的\Term{surjective}{满射性}得出。

\begin{SGAresult}{推论}{6.4}
在 6.2 的假设下，设 $n=\dim X$ 为偶数，或者 $\operatorname{car}(k)\neq2$，并设 $X$ 满足条件 (LV)。则存在整数 $d_0$，使得对 $d\geq d_0$，用\Term{degree}{次数}为 $d$ 的\Term{hypersurface}{超曲面}截 $X$ 所得的每个\Term{lefschetz-pencil}{Lefschetz 铅笔}都满足条件 (A)。
\end{SGAresult}
\SourcePage{331}
\begin{proof}
由 (6.3.2) 中的\Term{equivalence}{等价}性 $(\mathrm{i})\Longleftrightarrow(\mathrm{vi})$，只需证明：若用 $Y(d)$ 表示 $X$ 的一个\Term{degree}{次数}为 $d$ 的\Term{smooth}{光滑}\Term{hypersurface-section}{超曲面截面}，则
\[
\dim H^{n-1}(Y(d),\mathbb Q_\ell)\longrightarrow\infty
\quad\text{当 }d\longrightarrow\infty.
\tag{6.4.1}
\]
事实上，有如下精确结果：
\end{proof}

\begin{SGAresult}{\Term{lemma}{引理}}{6.4.2}
设 $X$ 是\Term{algebraically-closed-field}{代数闭域} $k$ 上的 $n\geq1$ 维\Term{smooth}{光滑}\Term{connected}{连通}\Term{projective}{射影}\Term{algebraic-variety}{簇}，用 $Y(d)$ 表示一个\Term{degree}{次数}为 $d$ 的\Term{hypersurface-section}{超曲面截面}。则存在\Term{polynomial}{多项式} $f(T)\in\mathbb Z[T]$，使得
\[
f(T)=\deg(X)T^n+\text{一个次数不超过 }n-1\text{ 的多项式},
\tag{6.4.2.1}
\]
\[
\text{对 }d\geq1,\qquad f(d)=\dim H^{n-1}(Y(d),\mathbb Q_\ell).
\tag{6.4.2.2}
\]
\end{SGAresult}
\begin{proof}
由第 XVII 讲 5.7.5，用 $C(X)$ 表示 $X$ 的\Term{total-chern-class}{总 Chern 类}，用 $H$ 表示一个\Term{hyperplane}{超平面}的\Term{first-chern-class}{第一 Chern 类}，则 $Y(d)$ 的\Term{euler-poincare-characteristic}{Euler--Poincaré 特征数}满足如下公式：
\[
\chi(Y(d))=\int_X\frac{dH\,C(X)}{1+dH}
=\sum_{i=1}^{n}(-1)^{i-1}d^i\int_X H^i\cdot C_{n-i}.
\tag{6.4.2.3}
\]
\SourcePage{332}
另一方面，由\Term{poincare-duality}{Poincaré 对偶性}及\Term{weak-lefschetz-theorem}{弱 Lefschetz 定理}，有
\[
\begin{split}
\chi(Y(d))={}&(-1)^{n-1}\dim H^{n-1}(Y(d),\mathbb Q_\ell)\\
&+2\sum_{i\geq1}(-1)^{n-1-i}\dim H^{n-1-i}(X,\mathbb Q_\ell).
\end{split}
\tag{6.4.2.4}
\]
合并 (6.4.2.3) 与 (6.4.2.4)，便得到所需结论。
\end{proof}

\textbf{6.5. 回顾 \Term{picard-lefschetz-formula}{Picard--Lefschetz 公式}（第 XV 讲 3.4）。}

设 $S$ 是一个\Term{strictly-henselian}{严格 Hensel}\Term{valuation-ring}{赋值环}的谱，$\eta$ 是它的\Term{generic-point}{泛点}，$\bar\eta$ 是 $\eta$ 上方的一个\Term{geometric-point}{几何点}，$s$ 是它的\Term{closed-point}{闭点}。设
\[
p:X\longrightarrow S
\]
是一个\Term{proper}{固有}\Term{morphism}{态射}，其中 $X$ 是 $n$ 维\Term{regular}{正则}\Term{connected}{连通}\Term{scheme}{概形}；$p$ 在一个点 $x$ 以外处处\Term{smooth}{光滑}，点 $x$ 位于\Term{special-fibre}{特殊纤维}上，且为\Term{special-fibre}{特殊纤维}的一个\Term{ordinary-quadratic-singularity}{普通二次奇点}。

用 $I$ 表示 $S$ 在 $s$ 处的\Term{inertia-group}{惯性群}。已知（第 XV 讲 3.4），当 $i\neq n-1$ 时，$I$ 在 $H^i(X_{\bar\eta},\mathbb Q_\ell)$ 上平凡地\Term{action}{作用}，而它在 $H^{n-1}(X_{\bar\eta},\mathbb Q_\ell)$ 上的\Term{action}{作用}由以下公式给出：

\textbf{6.5.1.} 若 $n=2m+1$，则 $I$ 通过一个\Term{isomorphism}{同构}于 $\mathbb Z/2\mathbb Z$ 的\Term{quotient-group}{商群}\Term{action}{作用}（注意：若 $\operatorname{car}k\neq2$，则这样的\Term{quotient-group}{商群}只有一个）。存在元素 $\delta\in H^{n-1}(X_{\bar\eta},\mathbb Q_\ell(m))$，它在相差一个\Term{signature}{符号}的意义下唯一确定，称为 $s$ 点处的\Term{vanishing-cycles}{消失圈}，并满足
\SourcePage{333}
\[
\delta\cdot\delta=(-1)^m2
\quad\text{于 }H^{2n-2}(X_{\bar\eta},\mathbb Q_\ell(n-1))\simeq\mathbb Q_\ell\text{ 中}.
\tag{6.5.1.1}
\]
\textbf{(6.5.2.2)} $\mathbb Z/2\mathbb Z$ 的非平凡元素 $\sigma$ 以如下\Term{reflection}{反射}\Term{action}{作用}于 $H^i(X_{\bar\eta},\mathbb Q_\ell)$：
\[
\sigma x=x-(-1)^m(x\delta)\delta.
\]
［元素 $x\cdot\delta$ 位于 $H^{2n-2}(X_{\bar\eta},\mathbb Q_\ell(m))\simeq\mathbb Q_\ell(-m)$ 中，故乘积 $(x\cdot\delta)\delta$ 定义 $H^{n-1}(X_{\bar\eta},\mathbb Q_\ell)$ 中的一个元素。］

\textbf{6.5.2.} 若 $n=2m+2$，则 $I$ 通过其\Term{canonical}{典范}\Term{quotient-group}{商群} $\mathbb Z_\ell(1)$ \Term{action}{作用}。存在元素 $\delta\in H^{n-1}(X_{\bar\eta},\mathbb Q_\ell(m))$，它在相差一个\Term{signature}{符号}的意义下\Term{canonical}{典范}确定，称为 $s$ 点处的\Term{vanishing-cycles}{消失圈}，使得 $\sigma\in\mathbb Z_\ell(1)$ 以如下\Term{transvection}{横切变换（transvection）}\Term{action}{作用}于 $x\in H^{n-1}(X_{\bar\eta},\mathbb Q_\ell)$：
\[
\sigma(x)=x-(-1)^m\sigma\cdot(x\cdot\delta)\delta.
\tag{6.5.2.1}
\]
［乘积 $x\cdot\delta$ 位于 $H^{2n-2}(X_{\bar\eta},\mathbb Q_\ell(m))\simeq\mathbb Q_\ell(-m-1)$ 中，因此将它乘以 $\sigma\in\mathbb Z_\ell(1)$ 得到元素 $\sigma\cdot(x\cdot\delta)\in\mathbb Q_\ell(-m)$；再将此元素乘以 $\delta\in H^{n-1}(X_{\bar\eta},\mathbb Q_\ell(m))$，便得到元素 $\sigma\cdot(x\cdot\delta)\delta\in H^{n-1}(X_{\bar\eta},\mathbb Q_\ell)$。］

再回到整体情形：

\begin{SGAresult}{\Term{theorem}{定理}}{6.6}
设 $X$ 是\Term{algebraically-closed-field}{代数闭域} $k$ 上的 $n$ 维\Term{smooth}{光滑}\Term{connected}{连通}\Term{projective}{射影}\Term{algebraic-variety}{簇}。假设 $n$ 为偶数或 $\operatorname{car}k\neq2$。设 $D$ 是 $X$ 的一个\Term{lefschetz-pencil}{Lefschetz 铅笔}，$\eta$ 是它的\Term{generic-point}{泛点}，$\bar\eta$ 是 $\eta$ 上方的一个\Term{geometric-point}{几何点}，
\SourcePage{334}
$s_1,\ldots,s_N$ 是 $D\cap\check X$ 中的各点。对每个 $i$，选取 $\bar\eta$ 与 $S_i$ 的\Term{geometric-generic-point}{几何泛点}之间的一个\Term{isomorphism}{同构} $\ell_i$，其中 $S_i$ 是 $D$ 在 $s_i$ 处的\Term{strict-henselization}{严格 Hensel 化}。这样，$S_i$ 在 $s_i$ 处的\Term{inertia-group}{惯性群} $I_i$ \Term{action}{作用}于 $H^{n-1}(X_{\bar\eta},\mathbb Q_\ell(j))$，并定义 $s_i$ 处的\Term{vanishing-cycles}{消失圈}（6.5）
\[
\delta_i\in H^{n-1}(X,\mathbb Q_\ell(m)),\qquad
m=\begin{cases}
\dfrac{n-1}{2},&n\text{ 为奇数},\\[4pt]
\dfrac n2-1,&n\text{ 为偶数}.
\end{cases}
\]
设 $\ell_i$ 的选择使各个 $I_i\subset\pi=\pi_1^{\mathrm{mod}}(D-D\cap\check X,\bar\eta)$ 在\Term{topology}{拓扑}意义下生成群 $\pi_1$（由 SGA 1，第 XIII 讲 2.10，这是可能的）。则有：

\textbf{6.6.1.} 若 $X$ 满足条件 (LV)，则\Term{vanishing-cohomology}{消失上同调}\Term{space}{空间} $E^{n-1}(X_{\bar\eta},\mathbb Q_\ell(m))$（定义为 $H^{n-1}(X_{\bar\eta},\mathbb Q_\ell(m))$ 中 $H^{n-1}(X,\mathbb Q_\ell(m))$ 的\Term{orthogonal-complement}{正交补}，参见 5.4.4）由\Term{vanishing-cycles}{消失圈} $\delta_i$，$i=1,\ldots,N$ 生成。

\textbf{6.6.2.} 各个\Term{vanishing-cycles}{消失圈} $\delta_i$ 在相差一个\Term{signature}{符号}的意义下，在整体\Term{monodromy-group}{单值群} $\pi_1(D-D\cap\check X)$ 的\Term{action}{作用}下彼此\Term{conjugacy}{共轭}。
\end{SGAresult}
\begin{proof}
为证明 6.6.1，注意到，由\Term{picard-lefschetz-formula}{Picard--Lefschetz 公式} 6.5，记 $\pi=\pi_1(D-D\cap\check X)$，并利用各 $I_i$ 生成 $\pi$ 这一事实，有
\SourcePage{335}
\[
\begin{aligned}
H^{n-1}(X_{\bar\eta},\mathbb Q_\ell(m))^\pi
&=\text{$H^{n-1}(X_{\bar\eta},\mathbb Q_\ell(m))$ 中}\\
&\qquad\text{由各 $\delta_i$ 生成的子空间的正交补}.
\end{aligned}
\tag{6.6.3}
\]
由 5.6.10，有
\[
H^{n-1}(X_{\bar\eta},\mathbb Q_\ell(m))^\pi
=\operatorname{Im}\bigl(H^{n-1}(X,\mathbb Q_\ell(i))\bigr),
\tag{6.6.4}
\]
故对 (6.6.3) 两边取\Term{orthogonal-complement}{正交补}，即得 6.6.1。

为证明 6.6.2，只需注意到，\Term{vanishing-cycles}{消失圈} $\delta_i$ 由 $s_i$ 处\Term{inertia-group}{惯性群}的\Term{representation}{表示}确定到相差一个\Term{signature}{符号}。因此，只需应用 6.1.2：它给出如下复合\Term{inertia-homomorphism}{惯性同态}在 $\pi_1^{\mathrm{mod}}(\mathbb P^r-\check X,\bar\eta)$ 中的\Term{conjugacy}{共轭}性：
\[
\begin{tikzcd}[column sep=large]
\mathbb Z_\ell(1)\arrow[r]&\pi_1^{\mathrm{mod}}(D-D\cap\check X,\bar\eta)\arrow[d]\\
&\pi_1^{\mathrm{mod}}(\mathbb P^r-\check X,\bar\eta).
\end{tikzcd}
\]
于是，它们所确定的\Term{vanishing-cycles}{消失圈}在 $\pi_1^{\mathrm{mod}}(\mathbb P^r-\check X,\bar\eta)$ 中元素的\Term{action}{作用}下，在相差一个\Term{signature}{符号}的意义下彼此\Term{conjugacy}{共轭}；再由 6.1.1 给出的\Term{surjective}{满射性}
\[
\pi_1^{\mathrm{mod}}(D-D\cap\check X)\longrightarrow\pi_1^{\mathrm{mod}}(\mathbb P^r-\check X),
\]
证明即告完成。
\end{proof}

\SourcePage{336}
\begin{SGAresult}{推论}{6.7}
在 6.6 的假设下，设 $X$ 满足条件 (LV)。则有：

\textbf{6.7.1.} $\pi_1(D-D\cap\check X)$ 在 $E^{n-1}(X_{\bar\eta},\mathbb Q_\ell)$ 上的\Term{representation}{表示}是\Term{absolutely-irreducible}{绝对不可约}的。

\textbf{6.7.2.} 若 $n$ 为偶数，则 $\pi_1(D-D\cap\check X)$ 的每个\Term{finite-index-subgroup}{有限指数子群}都在 $E^{n-1}(X_{\bar\eta},\mathbb Q_\ell)$ 上\Term{absolutely-irreducible}{绝对不可约}地\Term{action}{作用}。
\end{SGAresult}
\begin{proof}
显然，可以将所考虑的\Term{module}{模}替换为其\Term{twist}{扭曲} $E=E^{n-1}(X_{\bar\eta},\mathbb Q_\ell(m))$。设 $\bar E=E\otimes_{\mathbb Q_\ell}\overline{\mathbb Q}_\ell$，并取 $\bar E$ 中的非零元素 $x$。为证明 6.7.1，需要证明 $x$ 及其在 $\pi$ 下的所有\Term{conjugacy}{共轭}元所生成的 $\overline{\mathbb Q}_\ell$ \Term{vector-subspace}{向量子空间}等于整个 $\bar E$。由 6.6.1 及 6.6.2，只需证明此\Term{submodule}{子模}包含某个\Term{vanishing-cycles}{消失圈} $\delta_i$。但由于 $X$ 满足 (LV)，\Term{cup-product}{杯积}\Term{form}{形式}在
\[
H^{n-1}(X,\mathbb Q_\ell(m))\hookrightarrow H^{n-1}(X_{\bar\eta},\mathbb Q_\ell(m))
\]
上的\Term{restriction}{限制}是\Term{nondegenerate}{非退化}的，因此它在\Term{orthogonal-complement}{正交补} $E^{n-1}(X_{\bar\eta},\mathbb Q_\ell(m))$ 上的\Term{restriction}{限制}也是\Term{nondegenerate}{非退化}的；于是由 (6.6.1)，至少有一个 $i$ 使得
\[
x\cdot\delta_i\neq0.
\]
对于这个 $i$，由 6.5，$s_i$ 处的\Term{inertia-group}{惯性群}中存在元素 $\sigma_i$，使得
\[
\sigma_i(x)-x=\delta_i\text{ 的一个非零倍数},
\tag{6.7.3}
\]
这就完成了 6.7.1 的证明。

\SourcePage{337}
至于 6.7.2，此时 $n-1$ 为奇数，故 $\delta\cdot\delta=0$，并且对 $\sigma_i$ 的每个 $k$ 次幂，反复应用 6.5.2 得到
\[
\sigma_i^k(x)=x+k\,\sigma_i(x\cdot\delta_i)\delta_i.
\tag{6.7.4}
\]
所以 $\bar E$ 的一个 $\mathbb Q_\ell$ \Term{submodule}{子模} $M$ 在 $\pi$ 的\Term{action}{作用}下\Term{stable}{稳定}，当且仅当它在一个\Term{finite-index-subgroup}{有限指数子群}的\Term{action}{作用}下\Term{stable}{稳定}，因为这两个条件均\Term{equivalence}{等价}于
\[
x\in M\quad\Longrightarrow\quad
\text{对每个满足杯积 }x\cdot\delta_i\neq0\text{ 的 }i,\ \delta_i\in M.
\]
于是归结为已经证明的 6.7.1。
\end{proof}

\SGAsection{}{参考文献}
\begin{enumerate}
\item[{[1]}] P. Deligne，\emph{Théorème de Lefschetz et critères de dégénérescence de suites spectrales}（Lefschetz \Term{theorem}{定理}与\Term{spectral-sequence}{谱序列}\Term{degeneration}{退化}判据），Pub. Math. IHÉS \textbf{35}（1968），259--278。
\item[{[2]}] S. Kleiman，\emph{Algebraic cycles and the Weil conjectures}（\Term{algebraic-cycle}{代数圈}与 Weil 猜想），载 \emph{Dix Exposés sur la cohomologie des schémas}，North Holland，1968。
\item[{[3]}] S. Lefschetz，\emph{L'analyse situs et la Géométrie Algébrique}（位置分析与\Term{algebraic-geometry}{代数几何}），Gauthiers-Villars，巴黎，1924。
\item[{[4]}] D. Lieberman，\emph{Higher Picard Varieties}（高阶 Picard 簇），Amer. J. Math.，\textbf{15}（4）（1968），1165--1199。
\item[{[5]}] D. Mumford，未发表。
\item[{[6]}] A. Weil，\emph{Variétés Kählériennes}（Kähler 簇），Hermann，巴黎，1958。
\end{enumerate}
